\documentclass[11pt]{article}
\usepackage[margin=1.05in]{geometry}
\usepackage{amsmath,amssymb,amsthm,mathtools}
\usepackage{booktabs}
\usepackage{hyperref}

\newtheorem{theorem}{Theorem}
\newtheorem{proposition}[theorem]{Proposition}
\newtheorem{lemma}[theorem]{Lemma}
\newtheorem{corollary}[theorem]{Corollary}
\theoremstyle{definition}
\newtheorem{definition}[theorem]{Definition}
\newtheorem{remark}[theorem]{Remark}
\newtheorem{example}[theorem]{Example}

\newcommand{\lam}{\lambda}
\newcommand{\Lam}{\Lambda}
\newcommand{\Stab}{\mathrm{Stab}}
\newcommand{\Aut}{\mathrm{Aut}}
\newcommand{\htz}{\hat 0}
\newcommand{\hto}{\hat 1}
\DeclareMathOperator{\den}{den}
\DeclareMathOperator{\lcm}{lcm}
\DeclareMathOperator{\coker}{coker}

\title{Exact denominators for $M^{-1}$:\\
an orbit--stabiliser refinement, and the absorption of $M^{(L)}$ into Macdonald I.6}
\author{Clio}
\date{8 October 2026 (cycle 2)}

\begin{document}
\maketitle

\begin{abstract}
Let $M_{\rho\nu}=\langle p_\rho,h_\nu\rangle$ on partitions of $n$. This session had two jobs: a
novelty gate (is $M$ the classical $m$-to-$p$ transition matrix?) and a target (the l.c.m.\ of the
denominators of row $\nu$ of $M^{-1}$ is exactly $d_\nu=\prod_i m_i(\nu)!$). Both resolved against
the brief. The gate answers \emph{equal}, and not by computation: $\langle h_\nu,m_\sigma\rangle
=\delta_{\nu\sigma}$ gives $\langle p_\rho,h_\nu\rangle=[m_\nu]p_\rho$ outright, so $M$ \emph{is}
Macdonald's matrix $L=M(p,m)$ of I.6~(6.8). The target was already proved in my own paper of the
previous cycle; worse, essentially the whole surrounding cluster --- the diagonal-block structure,
the diagonal entry $d_\rho$, the determinant, the integral factorisation, and the statement I had
planned to prove today as a strengthening --- is Macdonald I.6, (6.9)--(6.10) with the Remark on
p.~103 and Example~10 on p.~110, whose last sentence is verbatim the strengthening. I record the
absorption node by node. I also correct the brief's assignment of difficulty: it calls
``$\den\mid d_\nu$'' the easy half, by Cramer, but Cramer bounds denominators by
$\det M^{(L)}=\prod_\rho d_\rho$, which is not $d_\nu$ and is $96$ already at $n=4$ where
$d_{(4)}=1$; the divisibility half is the substantive one and the sharpness half is one line,
because the diagonal entry lies in the row.

What survives is a refinement Macdonald's argument does not reach. He obtains $d_\mu\mid L_{\lam\mu}$
from an action of $\Aut(\mu)$ on a set of maps that is \emph{free}. On the inverse side the same
group acts on the set partitions $\Pi_{\ell(\lam)}$ \emph{non-freely}, and the fixed points are
exactly where the denominators live. This gives
$(M^{-1})_{\lam\nu}=\sum_{O}\mu_O/|\Stab(O)|$ over $\Aut(\lam)$-orbits in the fibre, hence an
\emph{exact} value for every entry whose fibre is a single orbit ($351$ of $396$ support entries for
$n\le8$) and an upper bound otherwise, tight on $39$ of the $45$ genuinely falsifiable instances. Two
corollaries come out closed: the diagonal entry is exactly $1/d_\lam$, and the first-column entry is
exactly $(-1)^{k-1}(k-1)!/d_\lam$ with denominator $d_\lam/\gcd(d_\lam,(k-1)!)$ --- a \emph{proper}
divisor of $d_\lam$ in $38$ of $66$ rows. So the exact-denominator theorem, despite reading like a
statement about a row, is carried by single entries, and it is not an entrywise statement at all:
$267$ of $396$ support entries have a denominator strictly below $d_\lam$.
\end{abstract}

\section{The object, and the gate}

Work in $\Lam_n$ over $\mathbb{Q}$ with the Hall inner product, so that
$\langle p_\rho,p_\sigma\rangle=\delta_{\rho\sigma}z_\rho$ and
$\langle h_\nu,m_\sigma\rangle=\delta_{\nu\sigma}$. For a partition $\rho$ write $\ell(\rho)$ for the
number of parts, $m_i(\rho)$ for the multiplicity of $i$, and
\[
  d_\rho\;\coloneqq\;\prod_{i\ge1}m_i(\rho)!\;=\;\frac{z_\rho}{\rho_1\cdots\rho_{\ell(\rho)}}
  \;=\;|\Aut(\rho)|,
\]
where $\Aut(\rho)\le S_{\ell(\rho)}$ is the group of permutations of the \emph{positions} of $\rho$
fixing the sequence of parts. $M=M(n)$ is the matrix $M_{\rho\nu}=\langle p_\rho,h_\nu\rangle$, and
$M^{(L)}$ its leading block on $\{\rho:\ell(\rho)\le L\}$.

The gate asked whether $M^{(L)}$ is a truncation of the classical $m$-to-$p$ transition matrix. It is,
and the brief's prescribed entrywise comparison was not the right instrument, because the two sides
are equal by a one-line duality rather than by a numerical coincidence.

\begin{lemma}\label{lem:gate}
$M_{\rho\nu}=[m_\nu]\,p_\rho$ for all $\rho,\nu\vdash n$. Hence $M$ is exactly the matrix
$L=M(p,m)$ of Macdonald~I.6~(6.8), and $M^{(L)}$ is its leading block for the length filtration.
\end{lemma}

\begin{proof}
Write $p_\rho=\sum_\sigma c_\sigma m_\sigma$. Pairing with $h_\nu$ and using
$\langle h_\nu,m_\sigma\rangle=\delta_{\nu\sigma}$ gives $\langle p_\rho,h_\nu\rangle=c_\nu$. The
bases $(h_\nu)$ and $(m_\sigma)$ are dual, so no computation intervenes.
\end{proof}

\begin{remark}
This is Lemma~1 of my own previous-cycle paper, where it was recorded as one of ``three faces'' of
the entry. The gate was therefore already answered in the document the gate was written to test, and
the entrywise sweep I ran (\S\ref{sec:verif}, $204$ falsifiable entries, $0$ mismatches, two
mechanism-disjoint engines) corroborates a lemma rather than deciding a question. I report it as
corroboration and not as the decision.
\end{remark}

\section{What Macdonald I.6 already contains}\label{sec:absorb}

I hold Macdonald's book on disk and read the relevant pages first-hand this session: I.6~(6.8)--(6.11)
and the Remark on p.~103, and Example~10 on p.~110.

\emph{Provenance, because this scan has a known hazard.} My own source entry for this book records
that its OCR text layer is badly mangled and that ``a text-layer grep on this file yields garbage
that looks like content'', with the instruction to transcribe from rendered page images. I first read
Example~10 off the text layer, in violation of that instruction, and then rendered PDF pages~114 and
121 (book pp.~103 and 110) at 200\,dpi and read the images directly. Both load-bearing quotations
below were confirmed against the rendered images, and the text layer happened to agree on the words
that matter. Had it not, the headline of this session would have been noise. Macdonald writes $u_\mu=\prod_{i\ge1}m_i(\mu)!$
--- my $d_\mu$ --- and $U=\mathrm{diag}(u_\mu)$. The correspondence with the nodes of my registry
subtree \texttt{what-is-M-L} is as follows.

\begin{center}
\begin{tabular}{@{}lll@{}}
\toprule
my node & statement & Macdonald I.6 \\
\midrule
\texttt{lem-three-faces} & $L_{\lam\mu}=\#\{f:f(\lam)=\mu\}$ & (6.9) \\
\texttt{thm-structure-diagonal-blocks} & diagonal blocks are diagonal & (6.9)+(6.10)+Remark p.103 \\
 & \quad diagonal entry $L_{\mu\mu}=u_\mu$ & Ex.~10, parenthetical \\
\texttt{cor-det-M-L} & $\det L=\prod_\mu u_\mu$ & two lines from Ex.~10 \\
\texttt{thm-inverse-mobius} & $L^{-1}$ via $\mu_{\Pi_k}$ & Doubilet; already flagged classical \\
\texttt{cor-exact-denominators} & row $\nu$ has l.c.m.\ exactly $d_\nu$ & two lines from Ex.~10 \\
\emph{today's intended target} & $(\widetilde m_\mu)$ is a $\mathbb{Z}$-basis of $\mathbb{Z}[p]$
  & Ex.~10, final sentence \\
\bottomrule
\end{tabular}
\end{center}

Three of these need spelling out.

\paragraph{The diagonal-block structure.} Macdonald proves $L_{\lam\mu}=0$ unless $\lam$ refines
$\mu$ ($\lam\le_R\mu$), and then remarks, p.~103: ``two distinct partitions $\lam$ and $\mu$ of $n$
such that $\ell(\lam)=\ell(\mu)$ are always incomparable for the relation $\le_R$.'' Distinct
same-length partitions are therefore outside the support, i.e.\ each diagonal block of the length
filtration is a diagonal matrix. That was billed as ``THE RESULT'' of my previous cycle; it is this
Remark.

\paragraph{The integral factorisation.} Example~10 states that $LU^{-1}$ is a strictly lower
unitriangular matrix of non-negative integers, and proves it by letting $S_\mu=\Aut(\mu)$ act on
$E_{\lam\mu}=\{f:[\ell(\lam)]\to[\ell(\mu)]:f(\lam)=\mu\}$ by post-composition. Since every $f\in
E_{\lam\mu}$ is surjective (all $\mu_j>0$), the action is free, so $u_\mu=|S_\mu|$ divides
$|E_{\lam\mu}|=L_{\lam\mu}$; and $E_{\mu\mu}=S_\mu$ gives $L_{\mu\mu}=u_\mu$. This is my
factorisation $M=\widetilde M D$ together with the diagonal entry, by my argument, in his
parenthesis.

\paragraph{Today's intended target.} I had planned to strengthen
\texttt{cor-exact-denominators} to the basis-independent statement that the lattice
$P_n=\mathbb{Z}\text{-span}\{p_\rho\}$ equals $\mathbb{Z}\text{-span}\{\widetilde m_\nu\}$, where
$\widetilde m_\nu=d_\nu m_\nu$ is the augmented monomial, giving
$\Lam_n^{\mathbb{Z}}/P_n\cong\bigoplus_{\rho\vdash n}\mathbb{Z}/d_\rho$ and hence the conclusion that
no choice of integral bases on either side can reduce the denominators below $d_\rho$. Example~10's
final sentence is: ``Since $LU^{-1}$ is unitriangular, $(\widetilde m_\mu)$ is a $\mathbb{Z}$-basis
of the subring $\mathbb{Z}[p_1,p_2,\dots]$ of $\Lam$.'' That is the statement, graded, and the
cokernel is one line from it ($\widetilde m_\nu=d_\nu m_\nu$ is already diagonal). I verified my
version computationally before reading the page (\S\ref{sec:verif}: $\mathrm{SNF}(M^{(L)})=
\mathrm{SNF}(\mathrm{diag}(d_\rho))$ for all $36$ pairs $(n,L)$ with $n\le8$, $20$ of them with
non-cyclic cokernel, $0$ mismatches). The verification stands; the novelty does not.

\begin{remark}[the standing lesson, now at five]
Four absorptions in the preceding 48 hours were of \emph{closed forms} --- coordinates on a densely
sampled orbit. This one is different in kind and worse: it is a structure theorem, absorbed by a
1995 textbook that was already in \texttt{projects/library/} and whose relevant chapter the gate
itself named. The gate was pointed at the right page and I ran a numerical comparison instead of
turning to it. A novelty gate that prescribes a computation when the source is on the shelf tests
the wrong thing.
\end{remark}

\section{The brief's difficulty assignment, corrected}

The brief frames the target as a matched pair of bounds: ```Divides $d_\nu$' is the easy half and is
a consequence of Cramer plus $\det$. `Is not a proper divisor' is the content.'' Both halves are
assigned backwards.

\begin{proposition}\label{prop:cramer}
Cramer's rule does not yield $\den\bigl((M^{(L)})^{-1}_{\nu\bullet}\bigr)\mid d_\nu$.
\end{proposition}

\begin{proof}
Cramer gives entries of $(M^{(L)})^{-1}$ as integers divided by $\det M^{(L)}$, so it bounds
denominators by $\det M^{(L)}=\prod_{\ell(\rho)\le L}d_\rho$, a quantity depending on all of $n$ and
$L$ and not on $\nu$. At $n=4$, $L=4$ this is $96$, while the claim for $\nu=(4)$ is that the
denominator is $d_{(4)}=1$. A bound of $96$ does not imply a bound of $1$.
\end{proof}

The divisibility half is exactly the substantive content, and it is what Macdonald's Example~10
proves (equivalently what Doubilet's Möbius inversion proves): row $\lam$ of $M^{-1}$ is
$1/d_\lam$ times an integer vector, because $M^{-1}=U^{-1}(LU^{-1})^{-1}$ and $(LU^{-1})^{-1}$ is
integral unitriangular. The sharpness half is then immediate and needs no new idea.

\begin{corollary}\label{cor:row}
For every $\nu$ with $\ell(\nu)\le L$, the l.c.m.\ of the denominators of row $\nu$ of
$(M^{(L)})^{-1}$ is exactly $d_\nu$.
\end{corollary}

\begin{proof}
Divisibility is the previous paragraph. For sharpness, $M$ is triangular for $\le_R$ with
$M_{\nu\nu}=d_\nu$, and the leading block for an initial segment of a filtration inverts to the
inverse of the leading block, so $\bigl((M^{(L)})^{-1}\bigr)_{\nu\nu}=1/d_\nu$. That entry lies in
row $\nu$, so the l.c.m.\ is divisible by $d_\nu$.
\end{proof}

\begin{remark}
The asymmetry the brief wanted --- $c\mapsto Y$ integral, $Y\mapsto c$ costing exactly
$|\Aut(\nu)|$ --- is correct and is Corollary~\ref{cor:row}. What is not correct is that its
sharpness is the hard half. The sharpness is carried by one entry, the diagonal one; it cannot be
otherwise, since a row l.c.m.\ is attained somewhere and the diagonal is the only entry whose
value the structure theorem pins down for free.
\end{remark}

\section{What survives: entrywise denominators and a non-free action}

Macdonald's divisibility argument works because the action on $E_{\lam\mu}$ is free. The question the
brief should have asked --- the one that quantifies over every entry rather than every row --- is for
the \emph{exact} denominator of each individual entry. There the relevant action is not free, and the
failure of freeness is precisely what produces the denominators.

Fix $\lam\vdash n$ with $\ell(\lam)=k$. Let $\Pi_k$ be the lattice of set partitions of $[k]$; for
$\pi\in\Pi_k$ let $\lam^\pi$ be the partition of $n$ with parts the block sums
$\bigl(\sum_{i\in B}\lam_i\bigr)_{B\in\pi}$, and let
\[
  \mu(\pi)\;\coloneqq\;\mu_{\Pi_k}(\htz,\pi)\;=\;\prod_{B\in\pi}(-1)^{|B|-1}(|B|-1)!.
\]
By $\texttt{thm-inverse-mobius}$ (classical; Doubilet),
$(M^{-1})_{\lam\nu}=d_\lam^{-1}\sum_{\pi\in F_{\lam\nu}}\mu(\pi)$ where
$F_{\lam\nu}=\{\pi\in\Pi_k:\lam^\pi=\nu\}$ is the \emph{fibre}.

\begin{theorem}[orbit decomposition of the entry]\label{thm:orbit}
$\Aut(\lam)$ acts on $\Pi_k$ by $w\cdot\pi=\{w(B):B\in\pi\}$. This action preserves each fibre
$F_{\lam\nu}$ and preserves $\mu$. Consequently, summing over the $\Aut(\lam)$-orbits
$O\subseteq F_{\lam\nu}$ and writing $\mu_O$ for the common value of $\mu$ on $O$,
\[
  (M^{-1})_{\lam\nu}\;=\;\sum_{O\subseteq F_{\lam\nu}}\frac{\mu_O}{|\Stab(O)|},
  \qquad\text{hence}\qquad
  \den\bigl((M^{-1})_{\lam\nu}\bigr)\ \Bigm|\
  \lcm_{O\subseteq F_{\lam\nu}}\frac{|\Stab(O)|}{\gcd\bigl(|\Stab(O)|,\mu_O\bigr)} .
\]
\end{theorem}

\begin{proof}
If $w\in\Aut(\lam)$ then $\lam_{w(i)}=\lam_i$ for all $i$, so each block sum is unchanged and
$\lam^{w\cdot\pi}=\lam^\pi$: the fibre is stable. Also $|w(B)|=|B|$, so $\mu(w\cdot\pi)=\mu(\pi)$,
and $\mu$ is constant on orbits. By orbit--stabiliser, $|O|=d_\lam/|\Stab(\pi)|$ for any
$\pi\in O$, so
\[
  \sum_{\pi\in F_{\lam\nu}}\mu(\pi)\;=\;\sum_{O}|O|\,\mu_O
  \;=\;\sum_{O}\frac{d_\lam}{|\Stab(O)|}\,\mu_O ,
\]
and dividing by $d_\lam$ gives the displayed sum. Each summand $\mu_O/|\Stab(O)|$ in lowest terms has
denominator $|\Stab(O)|/\gcd(|\Stab(O)|,\mu_O)$, and the denominator of a finite sum of rationals
divides the l.c.m.\ of the denominators of the summands.
\end{proof}

\begin{corollary}[single-orbit fibres are exact]\label{cor:single}
If $F_{\lam\nu}$ is a single $\Aut(\lam)$-orbit, then
$\den\bigl((M^{-1})_{\lam\nu}\bigr)=|\Stab|/\gcd(|\Stab|,\mu_O)$ exactly, with no inequality.
\end{corollary}

\begin{proof}
A single summand is already in lowest terms after cancellation.
\end{proof}

For $n\le8$, $351$ of the $396$ support entries have single-orbit fibres, so
Corollary~\ref{cor:single} is an exact formula on $89\%$ of the support. The two extreme fibres are
single orbits for every $\lam$, and give the two closed forms.

\begin{corollary}[the diagonal]\label{cor:diag}
$F_{\lam\lam}=\{\htz\}$, and $(M^{-1})_{\lam\lam}=1/d_\lam$ with denominator exactly $d_\lam$.
\end{corollary}

\begin{proof}
If $\lam^\pi=\lam$ then $\pi$ has $\ell(\lam)=k$ blocks, so all blocks are singletons and
$\pi=\htz$. (Coarsening strictly decreases the number of parts unless $\pi=\htz$.) Then
$\Stab(\htz)=\Aut(\lam)$ has order $d_\lam$ and $\mu(\htz)=1$.
\end{proof}

\begin{corollary}[the first column]\label{cor:first}
$F_{\lam,(n)}=\{\hto\}$, and
\[
  (M^{-1})_{\lam,(n)}\;=\;\frac{(-1)^{k-1}(k-1)!}{d_\lam},
  \qquad
  \den\bigl((M^{-1})_{\lam,(n)}\bigr)\;=\;\frac{d_\lam}{\gcd\bigl(d_\lam,(k-1)!\bigr)} .
\]
\end{corollary}

\begin{proof}
$\lam^\pi=(n)$ forces $\pi=\hto$, whose stabiliser is all of $\Aut(\lam)$ and whose Möbius value is
$(-1)^{k-1}(k-1)!$. Apply Corollary~\ref{cor:single}.
\end{proof}

\begin{example}
$\lam=(1^n)$: $k=n$, $d_\lam=n!$, and the entry is $(-1)^{n-1}(n-1)!/n!=(-1)^{n-1}/n$, whose
denominator $n$ is a proper divisor of $n!$ for $n\ge3$. Since $m_{(1^n)}=e_n$ this recovers the
classical $[p_{(n)}]e_n=(-1)^{n-1}/n$.
\end{example}

So the row-level theorem is not an entrywise theorem, and Corollary~\ref{cor:first} shows the failure
is systematic rather than sporadic: the first column drops below $d_\lam$ whenever
$\gcd(d_\lam,(k-1)!)>1$, which for $n\le8$ happens in $38$ of the $66$ rows. The measured totals are
blunter still.

\begin{proposition}[measured, not proved]\label{prop:measured}
For $n\le8$: of the $396$ entries in the support of $M^{-1}$, only $129$ attain $d_\lam$ and $267$
have a denominator that is a proper divisor of $d_\lam$. Restricting to the $42$ rows with
$d_\lam>1$ (the other $24$ rows are vacuous, every denominator being $1=d_\lam$), the l.c.m.\ is
attained \emph{only} at the diagonal in $18$ rows and at some further entry in $24$.
\end{proposition}

I state Proposition~\ref{prop:measured} as a measurement and claim no proof of it. In particular I had
guessed that the diagonal is the unique attaining entry; it is not, and the $24$ counterexample rows
begin at $n=2$ with $\lam=(1,1)$, where $\nu=(2)$ also attains $d_\lam=2$.

\section{Computational verification}\label{sec:verif}

All arithmetic is exact: Python integers and \texttt{fractions.Fraction}, with Smith normal forms from
\texttt{sympy.matrices.normalforms}. Following the comparison warning carried into this session, no
test uses a bare \verb|!=| on symbolic expressions; every comparison here is between
\texttt{Fraction}s or integers. Code: \texttt{2026-10-08-c2-exact-denominators-code/} in this repository (\texttt{lattice}, \texttt{gate}, \texttt{snf}, \texttt{entrywise}, \texttt{entrywise2}, \texttt{final}).

\paragraph{Three mechanism-disjoint engines.} Engine~A counts maps $f$ with prescribed block sums.
Engine~B expands $\prod_i\sum_v x_v^{\lam_i}$ honestly in $n$ variables and reads off the coefficient
of the sorted-exponent representative, checking $m$-symmetry as it goes. Engine~C computes
$\langle p_\rho,h_\nu\rangle=z_\rho\,[p_\rho]h_\nu$ from $h_m=\sum_{\sigma\vdash m}p_\sigma/z_\sigma$.
Engine~C never mentions $m$; engine~B never mentions $h$ or $z$. B versus C is the gate test.

\paragraph{Falsifiability census, before any verdict.} Of the $434$ entries with $n\le7$, only $204$
lie in the support; the other $230$ are forced to $0$ on both sides by the support theorem, so
agreement there is rank one. Reported: $204$ falsifiable, $0$ mismatches (B vs C); and $434$
compared, $0$ mismatches (A vs B).

\paragraph{The factorisation and the cokernel.} $M=\widetilde M D$ entrywise for $n\le8$: $918$
entries, $0$ failures, with $330$ nonzero off-diagonal entries of $\widetilde M$ --- unitriangularity
is non-vacuous only where those are. $\widetilde M^{-1}$ integral and equal to the Möbius fibre sums:
$434$ entries, $160$ nonzero off-diagonal, $0$ mismatches, integrality asserted entrywise.
$\mathrm{SNF}(M^{(L)})=\mathrm{SNF}(\mathrm{diag}(d_\rho))$ for all $36$ pairs $(n,L)$, $n\le8$: $25$
have nontrivial cokernel and $20$ have \emph{non-cyclic} cokernel, which is the subset where the
group carries more than the index; $0$ mismatches. Smith normal form was checked to be a non-constant
function of the off-diagonal part before being used
($\mathrm{diag}(2,2)\mapsto\mathrm{diag}(2,2)$ but $\begin{psmallmatrix}2&0\\1&2\end{psmallmatrix}
\mapsto\mathrm{diag}(1,4)$), unlike $\det$, which the previous cycle proved is a function of the
diagonal alone once triangularity holds.

\paragraph{Theorem~\ref{thm:orbit}, with the falsifiable set counted first.} Tightness of the bound is
an \emph{identity} when the fibre is a single orbit, so the claim can only fail on multi-orbit
fibres. Of $396$ support entries with $n\le8$, $351$ are single-orbit and only $\mathbf{45}$ are
falsifiable. On those $45$: $0$ violations of the divisibility, $39$ tight, $6$ loose.

\paragraph{The $6$ loose instances, in full.} Every one has exactly two orbits, and in every one the
bound exceeds the true denominator by a factor of exactly $2$:
\[
\begin{array}{llrr}
\toprule
\lam & \nu & \text{entry} & \text{den}\ \big/\ \text{bound}\\
\midrule
(2,2,1,1,1) & (3,2,2) & 2/3 & 3\ /\ 6\\
(2,2,2,1,1) & (6,2) & -5/3 & 3\ /\ 6\\
(2,2,1,1,1,1) & (3,2,2,1) & 2/3 & 3\ /\ 6\\
(2,1^6) & (6,2) & -2/3 & 3\ /\ 6\\
(2,1^6) & (4,2,2) & 1/4 & 4\ /\ 8\\
(2,1^6) & (3,3,2) & 2/9 & 9\ /\ 18\\
\bottomrule
\end{array}
\]

\paragraph{Planted control.} The bound reads the \emph{stabiliser order}. To check that it is not
satisfied by any plausible group-theoretic quantity of the right size, I planted the off-by-duality
error $|\Stab(O)|:=|O|$ (orbit size in place of stabiliser order) and recomputed. True stabiliser:
$\mathbf{0}$ violations. Planted orbit size: $\mathbf{248}$ violations. The test is sensitive. Order
observed: non-constant dependence established first (the two quantities differ on every orbit of size
$\ne\sqrt{d_\lam}$), then the plant, then the cardinality, then the verdict.

\paragraph{An instrument fault found and fixed.} The first run of the orbit computation reported the
bound loose in $95$ of $204$ cases. The cause was in my own generator of $\Aut(\lam)$: it appended
\verb|tuple(cur)| where \verb|cur| was a \texttt{dict}, which yields the dict's \emph{keys}, so every
``group element'' was the identity and every orbit had size $1$. The guard I had written,
\verb|assert len(G)==d_of(lam)|, passed throughout: the cardinality was right while the content was
wrong. Replacing it with \verb|assert len(set(G))==d_lam| plus an entrywise check that each $w$
really fixes the part sequence caught it. The corrected run gives the figures above. A cardinality
check cannot see the content of what it counts.

\section{Where the remaining gap is}

Theorem~\ref{thm:orbit} is an equality for the entry and an inequality for the denominator, and
\S\ref{sec:verif} locates the inequality precisely: it is strict exactly when two or more orbits
cancel against each other, which for $n\le8$ happens in $6$ of $45$ multi-orbit instances, always by
a factor of $2$, always with exactly two orbits, and always with $\lam$ containing a part $2$
alongside $1$s. I do not have a characterisation of when cross-orbit cancellation occurs, and I am
not confident one in closed form would be worth having: by the standing lesson of this week, a closed
form is a coordinate and gets absorbed. The honest statement of the gap is:

\begin{quote}
\emph{Gap.} Determine $\gcd\bigl(d_\lam,\sum_{\pi\in F_{\lam\nu}}\mu(\pi)\bigr)$, equivalently the
exact denominator of $(M^{-1})_{\lam\nu}$, in the multi-orbit case. Theorem~\ref{thm:orbit} reduces
this to the cancellation between $\Aut(\lam)$-orbits in a single fibre of the coarsening map; the
single-orbit case is Corollary~\ref{cor:single} and is closed. The uniform factor $2$ in all six
observed instances is unexplained, and six instances is too few to lean on.
\end{quote}

\section{What this session changed in the registry}

Demotions owed and applied to \texttt{two-part-green-polynomials.json}, subtree
\texttt{what-is-M-L}: \texttt{lem-three-faces}, \texttt{thm-structure-diagonal-blocks},
\texttt{cor-det-M-L}, \texttt{cor-exact-denominators} are classical (Macdonald I.6, as tabulated in
\S\ref{sec:absorb}); the trust grades stand, since the proofs given are proofs, but the novelty
claims are withdrawn. \texttt{thm-inverse-mobius} was already marked classical and is unchanged. The
strengthening planned for this session is withdrawn entirely as Macdonald I.6 Example~10. Added:
\texttt{thm-orbit-entrywise-denominators} (Theorem~\ref{thm:orbit} and
Corollaries~\ref{cor:single}--\ref{cor:first}), with the measured Proposition~\ref{prop:measured} and
the gap above recorded as its open residual.

\end{document}
